\chapter{Panoramica del sistema}
\label{ch:overview}

\section{Obiettivo del progetto}
FPGA-Neural implementa un \textbf{Neural Network Engine riusabile in hardware FPGA}.
L'insieme è composto da tre elementi: l'FPGA, che è il vero acceleratore; una RAM
dedicata fisicamente associata all'FPGA e non condivisa con l'host; e un'interfaccia
host indipendente dal sistema operativo, inizialmente SPI (con possibile estensione
futura a Dual~SPI).

Il principio fondante è la separazione fra chi \emph{esegue} il calcolo e chi lo
\emph{usa}: il calcolo della rete neurale avviene interamente dentro l'FPGA, mentre
il sistema host fornisce solo configurazione, parametri di rete, dati di ingresso,
controllo e lettura dei risultati. L'host non fa parte del datapath computazionale.
Sistemi host possibili includono SoC Linux, sistemi tipo Raspberry~Pi, ESP32,
microcontrollori e PC di sviluppo: la stessa architettura di engine deve poter essere
usata in sistemi completamente diversi.

\begin{center}
\begin{tikzpicture}[font=\footnotesize,node distance=8mm]
  \node[fnblockD,minimum width=42mm,minimum height=20mm] (host){\textbf{HOST}\\[2pt]
        {\scriptsize Configurazione}\\{\scriptsize Addestramento}\\{\scriptsize Controllo}};
  \node[fnblockT,below=14mm of host,minimum width=42mm,minimum height=20mm] (fpga)
        {\textbf{FPGA}\\[2pt]{\scriptsize Neural Network Engine}\\{\scriptsize Compute / Control}};
  \node[fnblock,below=14mm of fpga,minimum width=42mm,minimum height=13mm] (ram)
        {\textbf{RAM dedicata}\\{\scriptsize pesi / bias / buffer}};
  \draw[fnbus] (host) -- node[fnlbl,right]{SPI / Dual SPI} (fpga);
  \draw[fnbus] (fpga) -- node[fnlbl,right]{bus parallelo} (ram);
\end{tikzpicture}
\end{center}

\section{Configurazione hardware contro configurazione di rete}
Il progetto distingue con precisione fra l'\textbf{architettura hardware}
dell'acceleratore e i \textbf{parametri della rete neurale}.

L'architettura fisica dell'engine è definita al momento della sintesi e
dell'implementazione dell'FPGA. I parametri hardware tipici sono \code{N\_INPUTS},
\code{N\_NEURONS}, \code{N\_LAYERS}, \code{PARALLEL}, \code{DATA\_WIDTH},
\code{ACC\_WIDTH}: sono parametri Verilog risolti in fase di sintesi e determinano il
datapath contenuto nel bitstream. I parametri della rete --- pesi, bias, parametri di
attivazione e di quantizzazione, costanti specifiche --- vengono invece caricati a
runtime attraverso l'interfaccia host e memorizzati nella RAM associata all'FPGA.

\begin{fnnote}[Principio architetturale centrale]
Una build fissa il \emph{soffitto} della macchina (numero massimo di layer, larghezza
massima, \code{PARALLEL}); l'host configura la rete \emph{reale} --- numero di layer,
larghezza ingressi/uscite per-layer, attivazione per-layer e parametri addestrati ---
interamente a runtime, via SPI, nella memoria locale dell'FPGA. Un solo bitstream
serve qualunque topologia fino a quel soffitto.
\end{fnnote}

\section{Boot e inizializzazione}
L'FPGA viene configurato all'accensione tramite il consueto meccanismo di
configurazione (caricamento del bitstream da flash SPI). Il bitstream definisce
l'architettura hardware dell'engine; l'host non costruisce dinamicamente il datapath
durante il funzionamento normale, ma configura i dati di rete su cui il datapath già
esistente opera.

\begin{center}
\begin{tikzpicture}[font=\scriptsize,node distance=4.5mm,start chain=going below,
                    every node/.style={on chain}]
  \node[fnblockA,minimum width=60mm](p){Power-on};
  \node[fnblock,minimum width=60mm]{Configurazione FPGA (bitstream da flash)};
  \node[fnblockT,minimum width=60mm]{Neural Network Engine disponibile};
  \node[fnblock,minimum width=60mm]{Inizializzazione host (SPI)};
  \node[fnblock,minimum width=60mm]{Caricamento parametri di rete / pesi / bias};
  \node[fnblockD,minimum width=60mm]{Engine pronto};
  \begin{scope}[every path/.style={fnarrow}]
   \foreach \a/\b in {1/2,2/3,3/4,4/5,5/6}{}
  \end{scope}
  \foreach \i [count=\j from 2] in {1,...,5}{
     \draw[fnarrow] (chain-\i) -- (chain-\j);}
\end{tikzpicture}
\end{center}

\section{Addestramento e inferenza}
Addestramento e inferenza sono concettualmente separati. La prima implementazione non
richiede che l'FPGA esegua l'addestramento: i pesi possono essere calcolati
esternamente (PC/Linux/altro host) e trasferiti via SPI nella RAM dell'FPGA, che poi
esegue l'inferenza. Questo riduce drasticamente la complessità dell'hardware iniziale,
senza precludere una futura implementazione di training assistito o interamente
hardware (Fase~8 della roadmap, cap.~\ref{ch:roadmap}). Durante l'inferenza l'host
fornisce solo i dati di ingresso e recupera il risultato, ottenendo calcolo
deterministico, carico ridotto sull'host, parallelismo hardware, latenza prevedibile e
indipendenza dall'architettura della CPU host.

\section{Filosofia di progetto e riuso}
Il progetto va inteso come una \emph{piattaforma di accelerazione neurale FPGA
riusabile} più che come una singola rete. L'applicazione determina dimensione degli
ingressi, topologia, numero di layer e neuroni, parallelismo, precisione numerica,
funzioni di attivazione, requisiti di memoria e prestazioni; il processo di
generazione hardware produce l'implementazione FPGA corrispondente. La stessa
architettura HDL rimane concettualmente invariata mentre i parametri di sintesi
generano implementazioni appropriate ai diversi target applicativi.
